\section{Evaluation Protocol}
\label{sec:protocol}

The protocol fixes, before any result is read, what is measured, on which data, with which
denominators and with which uncertainty. Choices that are still open carry a marker.

\subsection{Data}
\label{sec:data}

Table~\ref{tab:data} lists the image sets. Image lists with file hashes are released.

\emph{Generated images.} We use three sources.
(i)~A \emph{reproduction set} drawn at random, with a fixed seed, from the eight validation subsets
of GenImage \pcite{zhu2023genimage}. The paper drew 100 images per subset. We draw 125 for the test
split, so that it has 1000 images, and 25 more for the development split. The 800 images of the
original evaluation are not identified, so the set cannot be the same.
(ii)~Images that we generate with Stable Diffusion 1.5 and 2.1 from class-name prompts. Version
1.5 is among the generators of GenImage; version 2.1 is not, and it is also the model inside the
attack.
\todo{prompts, classes, resolution, sampler, number of steps and seeds of source (ii); the pilot
used ``a photo of a'' followed by the class name, 512 pixels, 30 steps and guidance scale 7.5}
(iii)~Synthbuster \pcite{bammey2023synthbuster}, which provides 1000 images from each of nine
generators and is paired with the 1000 photographs of RAISE-1k \pcite{dangnguyen2015raise}.

\emph{Real images.} Real images are needed for false-positive rates, for the AUC and for threshold
calibration. For source (i) they are taken from the real folders of the same GenImage subsets. For
source (ii) they are ImageNet validation images \pcite{deng2009imagenet}, and the two sets of
source (ii) share them. For source (iii) they are the RAISE-1k photographs. The pilot runs used
Imagenette \pcite{howard2019imagenette}. The released code prepares no real images for the attack
experiments, and its test script and its training loader treat them in two other ways
(Section~\ref{sec:released-eval}). How real images are prepared in our experiments is therefore
our choice.

\begin{table}[!t]
\centering
\caption{Image sets and splits. ``dev.'': development split. The two Stable Diffusion sets share
their real images.}
\label{tab:data}
\small
\setlength{\tabcolsep}{4pt}
\begin{tabular}{@{}lrrrr@{}}
\toprule
 & \multicolumn{2}{c}{Generated} & \multicolumn{2}{c}{Real} \\
\cmidrule(lr){2-3}\cmidrule(l){4-5}
Set & test & dev. & test & dev. \\
\midrule
Reproduction set, 8 subsets  & $8\times125$ & $8\times25$  & 1000 & 2000 \\
Stable Diffusion 1.5         & 1000         & 200          & 1000 & 2000 \\
Stable Diffusion 2.1         & 1000         & 200          & \multicolumn{2}{c}{as above} \\
Synthbuster, 9 generators    & $9\times500$ & $9\times100$ & 500  & 500  \\
\bottomrule
\end{tabular}
\end{table}

\emph{Splits.} The development split serves everything that is tuned or calibrated: thresholds at
a target false-positive rate, the cap of the capped variant, the sensitivity sweep of
Section~\ref{sec:h1}, the admission of detectors and any hyper-parameter of a variant. The test
split is not used for any choice. The images and prompts of the pilot runs are kept out of all
splits: pilot files are excluded by hash, and generation uses other seeds.

\emph{Confirmatory set.} The decision rules of Section~\ref{sec:hypotheses} are applied to one set
only, the reproduction set. The other sets are analysed in the same way and reported without a
verdict.
\todo{GenImage has to be obtained first. If it cannot be, the Stable Diffusion 1.5 set becomes the
confirmatory set, the comparison with the published table is dropped, and both facts are stated as
deviations. Decide before the freeze.}

\subsection{Pre-processing and file format}
\label{sec:preprocessing}

\emph{Pre-processing.} Generated images are resized to $224\times224$ with a Lanczos filter,
without cropping, as the released driver does (\code{MAIN}~84). Real images are first cropped to
their central square and then resized with the same filter, as in the pilot, so that photographs
are not squashed; the generated images of our sources are square already. This is our main
pre-processing. It is not that of the release's test script, which resizes with a bilinear filter
without cropping, and not that of its training loader (Section~\ref{sec:released-eval}), and
resampling by itself moves detector scores
\pcite{rajan2025aligned,konstantinidou2026navigating}. Before any attack is run, we therefore
measure the clean operating point of each released detector under four pre-processings: the main
one; the main one without the central crop; that of the test script; and that of the training
loader with a central crop (Section~\ref{sec:results-clean}).
\todo{the documented training command leaves the load and crop sizes open; state the sizes used for
the fourth}

\emph{File format.} In GenImage the real images are JPEG files and the generated images are PNG
files, and detectors trained on it react to this difference \pcite{grommelt2024fakeorjpeg}. Source
(ii) has the same asymmetry by construction, and the output of the attack is always a newly written
PNG file. A comparison of real with generated or attacked images on the files as distributed is
therefore confounded with the format. The channel setting ``JPEG 95'' of
Section~\ref{sec:channel} passes all three kinds of images through the same encoder as the last
step before detection. We call it the \emph{format-equalised reference}, following the
recommendation of \tcite{Zhou and Wang}{zhou2026audit} to re-encode both classes in the same
format.

\emph{Which analysis decides.} H1 compares attacked images with each other and is decided in the
main analysis: main pre-processing, lossless files. H2 and H3 compare attacked images with
thresholds that come from real images. Their verdicts are computed three times: in the main
analysis; with the format-equalised reference in the place of the lossless files; and with the
pre-processing of the test script for real and unmodified generated images. A verdict on H2 or H3
stands only if it is the same in all three. Otherwise the hypothesis is undecided, and the three
analyses are reported side by side.

\subsection{Detectors}
\label{sec:detectors}

\emph{Within-family detectors.} The four detectors released with the attack (\detE, \detR, \detD,
\detS) are used with the released weights, the released wrapper and the threshold $\tau=0$ on the
logit. We do not retrain them.

\emph{Out-of-family detectors.} We add detectors from other groups, chosen among UnivFD
\pcite{ojha2023univfd}, DRCT \pcite{chen2024drct}, NPR \pcite{tan2024npr}, AIDE
\pcite{yan2025aide}, Effort \pcite{yan2025effort} and Community Forensics
\pcite{park2025communityforensics}. UnivFD is among the five additional detectors of the original
paper. The other four of those five, DIRE, LGrad, GFF and RECCE, are included if their public
weights load in our environment. \todo{try them; state for each one that is left out why}
A detector is admitted for a generated set if its clean AUC on the development split of that set
is at least 0.8; detectors that fail are listed with their clean performance and left out of the
attack tables for that set. The same rule applies to the four released detectors on sets other
than the reproduction set. Each out-of-family detector is used with its official pre-processing
and default threshold, and we record whether this pre-processing leaves a $224\times224$ input
unchanged. Scores are oriented so that higher means ``real''.
\todo{final list, after one test image has been run through each candidate}

For every detector we record the SHA-256 hash of the weight file, the source of the code and its
commit.

\subsection{Harness}
\label{sec:harness}

\emph{Running the release.} The code needs changes to its environment before it runs elsewhere. It
sets a proxy address and a GPU index (\code{MAIN}~\mbox{4--6}; \code{OA}~7; \code{TV}~\mbox{3--4}) and assumes a
CUDA device (\code{MAIN}~137; \code{OA}~81; \code{DLA}~227). Its detector loaders fetch ImageNet
weights from the network before they load the released weights (\code{OA}~58, 74). The weight
files are distributed through a file-sharing link in the README.
\verify{whether the default model identifier \code{stabilityai/}\allowbreak\code{stable-diffusion-2-1-base}
still resolves; the pilot runs used a mirror, whose snapshot hash must be recorded}

\emph{Departures from the release.} No line of the official files that takes part in the
computation is changed. Our harness imports the official functions for pre-processing, inversion,
the diffusion step, the detector wrappers and the \CVAE classes. It departs from the released
driver in five ways. (i)~The per-image loop of Algorithm~\ref{alg:released} is re-implemented;
every configuration of Section~\ref{sec:attacks} will be a set of command-line switches.
(ii)~A small shim lets the same code run on devices other than CUDA. (iii)~The environment differs
from the pinned one: PyTorch~2.2.2, torchvision~0.17.2, transformers~4.30.2, NumPy~1.26.4 and
SciPy~1.11.4 on Python~3.10, where the repository pins 2.0.0, 0.15.1, 4.42.4, 1.24.4 and 1.9.1 and
names Python~3.8. (iv)~The diffusion checkpoint is loaded from a mirror. (v)~In the Colab workflow
of the pilot, eight lines in six official files that set proxies and device names were rewritten
in place; on the laptop no file was changed. \todo{replace this patch by settings outside the
official files, so that the statement ``no official file is modified'' holds for every run}
\todo{Equivalence check, to be reported here: run the unmodified \code{main.py} and the harness
with default switches on the same images and the same device, and compare the saved files pixel by
pixel and the logged decisions one by one. Until this is done, statements about ``the released
pipeline'' in our experiments rest on our reading of the code and on the stand-in execution of
Appendix~\ref{app:verification}.}

\emph{One pass, shared intermediates.} With $\alpha=\epsilon$, a difference in the sign of one
gradient component changes a pixel by $4/255$ or $8/255$, so two separate runs of the same
configuration need not agree exactly. The harness therefore computes all configurations that share
an intermediate result in one pass per image and surrogate, from the same $\xone$, the same
$\xtwo$ and the same latents. For each image it stores, as 8-bit PNG files written with the
truncation of the released code, the input, $\xone$, $\xtwo$ and the output of every
configuration. It stores the scores of all detectors on each of them twice: on the floating-point
tensor, as the released driver logs them, and on the file read back from disk. Unless stated
otherwise, our measures use the files.
\todo{the pilot harness reaches the configurations \cfg{A2}, \cfg{A3}, \cfg{A5}, \cfg{X5} and
\cfg{X6}; it stores the input, the output of the plain decoder and the final output, and logs
tensor scores at four stages. Still to be added: storing $\xone$ and $\xtwo$, and the switches for
the control image, for the image that is encoded and for the output decoder}

\emph{Noise floor.} The released configuration is run twice on the same device and once on a
second device, and the share of decisions that differ between these runs is reported next to every
contrast between configurations. \todo{devices; the two pilot runs were made on different devices}

\emph{Safeguards.} Every prediction that the attack script records is preceded by an assertion
that the detector is in evaluation mode, and the scores of the unmodified images are compared
across runs before any other number is read. Appendix~\ref{app:bn} describes the fault that led to
these two rules.

\subsection{Pipeline configurations}
\label{sec:attacks}

Table~\ref{tab:configs} lists the configurations. Each is obtained from
Algorithm~\ref{alg:released} by choosing which image is encoded into the latent, whether and where
the loss of the latent optimisation is computed, which decoder produces the output, and which image
is the control image. The surrogate is part of a configuration: on the confirmatory set every
configuration is run with each of the four released detectors as surrogate.

Configurations \cfg{A1} to \cfg{A5} correspond to the rows of the published ablation
(Table~\ref{tab:paper-ablation}). Row 4 of that table, ``no PGD'', can be read in two ways, because
the second PGD call sits inside the \CVAE branch of the code
(Section~\ref{sec:released-pipeline}): either both calls are removed (\cfg{A4a}) or only the
pre-processing call is (\cfg{A4b}). Under Hypothesis~1 the two readings give opposite results, so we
run both and set both against row 4. Configurations \cfg{X1} to \cfg{X6} are ours. Two
configurations have a switch in the released driver, \cfg{A2} and \cfg{A3}. The two switches
cannot be combined: with both, the driver stops with an error, because a mask that the plain
decode of the output uses is defined inside the optimisation loop (\code{DLA}~432, 526). \cfg{X6}
therefore has no counterpart in the release, and neither have \cfg{A1}, \cfg{A4a}, \cfg{A4b} and
\cfg{X1} to \cfg{X5}.

\begin{widetable}
\centering
\caption{Pipeline configurations. ``Released'' latent optimisation: lines 3 to 7 of
Algorithm~\ref{alg:released}, with the loss on the plain decoder. ``Encoded'': the image that is
encoded into the latent. ``Control'': the control image. A dash: the stage is absent.}
\label{tab:configs}
\small
\setlength{\tabcolsep}{4pt}
\begin{tabular}{@{}lllll>{\raggedright\arraybackslash}p{0.40\linewidth}@{}}
\toprule
 & Encoded & Latent optimisation & Decoder & Control & Purpose or counterpart \\
\midrule
\cfg{A1}  & --      & --                    & --    & --            & output is $\xone$; row 1 of Table~\ref{tab:paper-ablation} \\
\cfg{A2}  & $\xone$ & none                  & skip  & $\xtwo$       & row 2; released switch \cmdarg{iterations}~\code{0} \\
\cfg{A3}  & $\xone$ & released              & plain & --            & row 3; released switch \cmdarg{is\_encoder}~\code{0} \\
\cfg{A4a} & $\xin$  & released              & skip  & $\xin$        & row 4, both PGD calls removed \\
\cfg{A4b} & $\xin$  & released              & skip  & $\PGD(\xin)$  & row 4, only the pre-processing call removed \\
\cfg{A5}  & $\xone$ & released              & skip  & $\xtwo$       & row 5; the released pipeline \\
\midrule
\cfg{X1}  & --      & --                    & --    & --            & output is $\Quant(\xtwo)$: the control image of the release alone \\
\cfg{X2}  & $\xone$ & released              & skip  & $\xin$        & clean control image \\
\cfg{X3}  & $\xone$ & released              & skip  & $\xone$       & second PGD call removed \\
\cfg{X4}  & $\xin$  & none                  & skip  & $\xin$        & reconstruction floor: no attack at all \\
\cfg{X5}  & $\xone$ & loss on skip decoder  & skip  & $\xtwo$       & loss computed on the decoder that produces the output \\
\cfg{X6}  & $\xone$ & none                  & plain & --            & plain route without optimisation \\
\bottomrule
\end{tabular}
\end{widetable}

The plain decoder is always used with $\kv$ and the skip decoder with $\kc$, as in the release.
The contrasts that we use are these. \cfg{A5} against \cfg{A2}: what the latent optimisation adds
to the released pipeline. \cfg{X2} against \cfg{X4}: what the route through the latent carries when
the control image is clean. \cfg{A5} against \cfg{A1}: what the pipeline adds to one PGD call.
\cfg{A5} against \cfg{X3}: what the second PGD call adds. \cfg{A5} against \cfg{X1}: what the
passage through the skip decoder changes. \cfg{X5} against \cfg{A2}: whether the latent
optimisation reaches the saved image once its loss is computed on the skip decoder. \cfg{A3}
against \cfg{X6}: whether it works on its own route.

\emph{Variants.} Four modifications of the released pipeline are examined as exploratory material:
a surrogate ensemble of \detE, \detR and \detD, with \detS held out, which averages their logits in
both PGD calls and their losses in the latent optimisation \pcite{liu2017delving}; a capped
objective, described in Section~\ref{sec:h2}; an LPIPS term in the loss of the latent
optimisation; and input diversity \pcite{xie2019diversity} in the latent optimisation.
\todo{weight of the LPIPS term and parameters of the input diversity; the pilot used weight 1 and
probability 0.5} The last two change only the latent optimisation. Under Hypothesis~1 they should
change almost no decision on the saved image, which we check with the statistic of
Section~\ref{sec:h1}. Variants with internal randomness are run with the seeds 0, 1 and 2.

\emph{Baselines.} PGD from $\xin$ with $\epsilon=8/255$, step size $2/255$ and 10 steps, on the same
objective; a transfer baseline that combines momentum, input diversity and a surrogate ensemble
\pcite{dong2018momentum,xie2019diversity,liu2017delving} \todo{its parameters}; and FPBA
\pcite{diao2024fpba}, whose code is public. \todo{confirm that FPBA runs on our detectors; otherwise
state why it is omitted} Together with \cfg{A1} and \cfg{X1}, the first baseline shows how the
released pipeline compares with pixel-space attacks that use one PGD call, two PGD calls, and one
call with twice the bound.

All configurations use the same image lists.

\subsection{Channel}
\label{sec:channel}

The following settings are applied to saved $224\times224$ files, identically for real images,
unmodified generated images and attacked images.
\begin{itemize}
\item \emph{None}: the lossless file.
\item \emph{JPEG} with quality 95, 75 and 50, encoded with Pillow at its default chroma
sub-sampling.
\item \emph{Resizing} by a factor of 0.75 and 0.5 with a Lanczos filter, followed by bicubic
up-sampling back to $224\times224$.
\item \emph{Composite settings} that follow the medium and high degradation levels of
\tcite{Mavali et al.}{mavali2024fake}, for comparability: down-sampling to 75\% with JPEG quality 50;
and Gaussian blur, down-sampling to 50\%, JPEG quality 20 and Gaussian noise.
\verify{their exact parameters}
\end{itemize}
The pilot used JPEG 75 and resizing by 0.5. For every setting, the set of images that the detector
catches before the attack is recomputed under that setting (Section~\ref{sec:metrics}).

This channel has a limit that we state here and not only at the end. A platform processes an image
at its native size, and the detector's pre-processing comes afterwards. We process the
$224\times224$ file, because that is the only output the released pipeline produces. We also apply
no cropping. Our channel results therefore describe the attack as released at its working
resolution and are not a simulation of a platform.

\subsection{Measures}
\label{sec:metrics}

Let $\Xgen$ be the generated test images and $\Xreal$ the real test images of a set, $a$ a
configuration, with $a(v)$ its output for the input $v$, and $h$ a channel setting. Let $s_d$ be
the score of detector $d$ and $\tau_d$ its threshold.

\emph{Clean operating point.} The set of generated images that the detector catches before the
attack, and the false-positive rate on real images, both under the channel setting $h$, are
\begin{align}
  \Bset_d(h) &= \{v\in\Xgen : s_d(h(v))\le\tau_d\}, \label{eq:detected}\\
  \FPR_d(h) &= \frac{|\{v\in\Xreal : s_d(h(v))\le\tau_d\}|}{|\Xreal|}. \label{eq:fpr}
\end{align}
We report $|\Bset_d(h)|/|\Xgen|$ and $\FPR_d(h)$ for every detector and channel setting. We call
the images in $\Bset_d(h)$ the \emph{caught} images.

\emph{Conditional success rate.} Success is counted only on caught images:
\begin{equation}
  \ASR_d(a,h)=\frac{|\{v\in\Bset_d(h) : s_d(h(a(v)))>\tau_d\}|}{|\Bset_d(h)|}.
  \label{eq:asr}
\end{equation}
The denominator is printed next to every value. We omit $h$ for lossless files. Whenever we write
``success rate'' for our own results we mean this quantity, as a fraction; the rates quoted from
the paper are percentages with an unstated denominator. For the comparison with the release we
also report the quantity that the released driver logs: the share of all attacked images that are
classified as real, computed on the floating-point output.

\emph{Discordance.} For two configurations $a$ and $a'$, the discordance $\psi_d(a,a')$ is the
share of the caught images whose outputs under $a$ and under $a'$ receive different decisions from
detector $d$. The difference of the two success rates cannot exceed it in absolute value.

\emph{Separability.} The area under the ROC curve (AUC) between real images and attacked images is
taken over all generated images and not only the caught ones:
\begin{equation}
\begin{aligned}
  \AUC_d(a,h)={}&\Pr\bigl[s_d(h(v_r))>s_d(h(a(v_g)))\bigr]\\
  &+\tfrac12\Pr\bigl[s_d(h(v_r))=s_d(h(a(v_g)))\bigr],
\end{aligned}
  \label{eq:auc}
\end{equation}
with $v_r$ drawn from $\Xreal$ and $v_g$ from $\Xgen$ \pcite{hanley1982meaning,fawcett2006roc}. A
value near 1 means that attacked images score below real images, as unmodified generated images
do. A value near 0 means that they score above real images: the two sets are separable, in the
reversed direction. A value near 0.5 does not show that the sets cannot be told apart: if half of
the attacked images score above all real images and half below, the AUC is 0.5 and a two-sided test
separates the sets completely. The AUC is therefore reported as a description, and the decision
rules use the calibrated tests below.

\emph{Calibrated tests.} Let $n$ be the number of real development images, $q$ a target
false-positive rate and $k=\lfloor qn\rfloor$. The \emph{one-sided test} flags an image if its
score is at most the $k$-th smallest score of the real development images under the same channel
setting. The \emph{two-sided test} uses $k'=\lfloor qn/2\rfloor$: it flags an image if its score
is at most the $k'$-th smallest or at least the $k'$-th largest of these scores. We call the second
event \emph{overshoot}. We use $q=0.05$, and $q=0.01$ where 2000 real development images are
available \pcite{carlini2022membership}. With 2000 development images, the false-positive rate
that such a threshold realises on new real images lies between 0.041 and 0.060 for $q=0.05$ and
between 0.006 and 0.015 for $q=0.01$, with probability 0.95; with 500 it lies between 0.033 and
0.071 for $q=0.05$. We therefore print the false-positive rate measured on the real test images
next to every rate of a calibrated test.

\emph{Distortion.} PSNR, SSIM \pcite{wang2004ssim}, LPIPS with the VGG backbone
\pcite{zhang2018lpips} and the maximum absolute pixel difference, all between the saved output and
the $224\times224$ input, as means and as 5\%, 50\% and 95\% quantiles over images. The bound
$\epsilon$ is a parameter of one stage of the pipeline and not a bound on its output, so we report
measured distortion.

\emph{Cost.} Seconds per image on named hardware, and the number of forward and backward passes
through each detector.

\subsection{Statistical reporting}
\label{sec:statistics}

\emph{Error level.} All intervals have an error level of 0.05. A two-sided interval puts 0.025 on
each side, and a one-sided bound is taken at 0.025.

\emph{Single proportions.} Every rate is given with its count and its Wilson interval
\pcite{wilson1927probable,brown2001interval}. Where a decision rule bounds a share that may be
zero, we use the exact binomial bound of \tcite{Clopper and Pearson}{clopper1934use}, because its
coverage is guaranteed. We write $\underline{p}$ and $\overline{p}$ for the exact lower and upper
bounds of a share $p$.

\emph{Paired comparisons.} Two configurations are always compared on the same caught images. Let
$m$ be their number, $n_{10}$ the number on which the first configuration succeeds and the second
does not, and $n_{01}$ the number for the opposite case. With $p_{10}=n_{10}/m$ and
$p_{01}=n_{01}/m$, the difference of the two success rates is $p_{10}-p_{01}$ and the discordance
is $p_{10}+p_{01}$. A bootstrap interval is of no use when no image changes its decision, since
every resample then gives zero. We use exact bounds instead. The interval for the difference is
$[\underline{p}_{10}-\overline{p}_{01},\;\overline{p}_{10}-\underline{p}_{01}]$, with each of the
four bounds at a quarter of the error level; a one-sided bound for the difference uses two of
these bounds, each at half of its error level. The coverage is at least the nominal one, so the intervals are
conservative. We also report the exact test of \tcite{McNemar}{mcnemar1947note} on the discordant
pairs \pcite{dietterich1998approximate}.

\emph{Equivalence.} To show the absence of a relevant difference, a non-significant test is not
enough. We use inclusion in a margin $\delta$ that is stated in advance
\pcite{schuirmann1987comparison,lakens2017equivalence}: two configurations are declared equivalent
for a detector if the exact upper bound of their discordance is below $\delta$. This is stricter
than equivalence of the two rates.

\emph{Several cells.} If a rule needs every cell to pass, each cell is tested at the full error
level. If one cell is enough to trigger a rule, each cell is tested at the error level divided by
the number of chances, which the rule states.

\emph{Ratios and AUC.} Intervals come from a percentile bootstrap \pcite{efron1979bootstrap} with
10\,000 resamples of the test images. Real and generated images are resampled separately, all
detectors and surrogates of an image stay together, and thresholds calibrated on the development
split stay fixed. At an AUC of 0 or 1 this interval has zero width; we then say so and give no
interval.

\emph{Sources of variation.} The released pipeline has no random start in PGD and a fixed seed for
the VAE posterior, so for a given image and device it is deterministic up to numerical effects,
which the noise floor of Section~\ref{sec:harness} measures. The variation that matters is across
images. We therefore report every headline number on three thirds of the test split, cut in the
order of the released image list, in addition to the pooled value. For a variant with internal
randomness each seed is reported.

\subsection{Scale and order of the runs}
\label{sec:scale}

The configurations of Table~\ref{tab:configs} need three latent optimisations per image and
surrogate: one for \cfg{A3}, \cfg{A5}, \cfg{X2} and \cfg{X3}, which differ only in the final decode,
one for \cfg{A4a} and \cfg{A4b}, and one for \cfg{X5}. The remaining configurations need none. For
Table~\ref{tab:configs} on the test split of the confirmatory set this is 12\,000 optimisations
for four surrogates. The other sets, with the configurations \cfg{A1}, \cfg{A2}, \cfg{A5},
\cfg{X1}, \cfg{X2} and \cfg{X4} and the surrogates \detE and \detS, add 13\,000. The sensitivity
sweep adds 5600 and the variants about 11\,000. In the pilot, one image took about 5 seconds on
an NVIDIA A100 and about 11 seconds on an Apple laptop with the MPS back end, for the whole
released pipeline with one optimisation. \todo{confirm with a timing run; at the pilot speed the
plan is about 60 hours on one A100} The order of work is: the clean operating point; the
equivalence check and the noise floor; all configurations on the confirmatory set; then the other
sets; then the variants and baselines; then the channel settings, which need no further attack
runs.

\subsection{What every table states}
\label{sec:checklist}

Each results table in this paper states, in its caption or in a note: the image set and the number
of images; the detector and whether it was among the surrogates; the denominator of every rate; the
threshold; the channel setting; the false-positive rate on real images under the same setting; and
the file format. The hashes of the detector weights, the switches of every configuration and the
software and hardware are given once, in the released material. This is the list to which we hold
our own tables.
